\documentclass[11pt,a4paper]{article}

\usepackage[
    top=1cm,
    bottom=2.5cm,
    left=2.5cm,
    right=2.5cm
]{geometry}
\usepackage{booktabs}
\usepackage{array}
\usepackage{caption}

\title{
    \textbf{SODS Final Project - Group 20}\\[0.5em]
    \large
    \begin{tabular}{ccc}
        Name Surname sID & Name Surname sID & Name Surname sID
    \end{tabular}
}

\begin{document}

\date{}

\maketitle
\vspace{-2cm}
\section{Baseline Synthesis}

\begin{table}[h]
\centering
\caption{Baseline synthesis results}
\begin{tabular}{lc}
\toprule
\textbf{Metric} & \textbf{Value} \\
\midrule
Clock period (ns) & 7 \\
Area ($\mu m^2$) & 17708.964844 \\
Slack (ns) & 0.394614 \\
Leakage power (mW) & 4.961e-04 \\
Dynamic power (mW) & 1.285 \\
\bottomrule
\end{tabular}
\end{table}

\section{Power Optimization}

\subsection{Clock-Gating}

\begin{table}[h]
\centering
\caption{Clock-gating optimization results}
\begin{tabular}{lc}
\toprule
\textbf{Metric} & \textbf{Value} \\
\midrule
Area ($\mu m^2$) & 13916.756836 \\
Slack (ns) & 0.099250 \\
Leakage power (mW) & 4.711e-04 \\
Dynamic power (mW) & 0.7102 \\
\bottomrule
\end{tabular}
\end{table}

\paragraph{Discussion.} Clock gating is a technique for power saving, so its main effect will be on leakage and dynamic power. In this case, both those metrics have reduced, dynamic power by 0.57 mW (44 per cent) and leakage power by 0.25e-04 mW (5 per cent). By stopping the clock to unchanging sections of the design dynamic power is drastically reduced since we are no longer unnecessarily powering the clock. Leakage power also decreases by the fact that keeping transistors completely off saves power compared to transistors that are off but only between abundant toggling. It only saves a bit of power however, because the dynamic power is the main target of clock gating. 
The slack will be negatively affected by clock gating, because we are adding more hardware in the clock path (clock gates), and in this case the slack decreased by 0.295 ns (75 per cent).
% Discuss and motivate the impact of clock gating on PPA metrics.

\subsection{Clock-Gating + Multi-VT}

\begin{table}[h]
\centering
\caption{Clock-gating and Multi-VT optimization results}
\begin{tabular}{lc}
\toprule
\textbf{Metric} & \textbf{Value} \\
\midrule
Percentage of LVT cells (\%) &  \\
Percentage of SVT cells (\%) &  \\
Percentage of HVT cells (\%) &  \\
Area ($\mu m^2$) &  \\
Slack (ns) &  \\
Leakage power (mW) &  \\
Dynamic power (mW) &  \\
\bottomrule
\end{tabular}
\end{table}

\paragraph{Discussion.} ...
% Discuss and motivate the impact of Multi-VT synthesis on PPA metrics.

\end{document}